% !PRINTED 454
\weberpage{454}
%<<CONT
bleiben zur Erfüllung von (23) nur folgende zwei Möglichkeiten übrig:

1.\quad $D=-3, \ x=1, \ y=\pm 1.$

Da es für die Diskriminante $-3$ nur eine Formenklasse gibt, so können wir $A=1$, $B=1$, $C=1$ annehmen, d.\,h. für $\omega$ die imaginäre dritte Einheitswurzel $e^{\frac{2\pi i}{3}}$ setzen, und erhalten unter den Größen (15) drei äquivalente, indem wir $y=+1$, und $=-1$ annehmen: % [sic] printed second equality omits y

\weq{24}{p\,\omega, \quad \frac{\omega}{p}, \quad \frac{\omega+1}{p},}

wie auch aus den Gleichungen:
\[
\frac{\omega}{p} = \frac{-1}{p+p\omega}, \quad \frac{\omega+1}{p} = \frac{-1}{p\omega}
\]
erkannt wird, aus denen die Äquivalenz der drei Größen (24) evident ist.

2.\quad $D=-4, \ x=1, \ y=0.$

Wir können ebenso $B=0$, $A=C=1$, d.\,h. $\omega=i$ annehmen und finden die zwei äquivalenten Werte:

\weq{25}{p\,\omega, \quad \frac{\omega}{p},}

wie auch aus der Gleichung:
\[
p\,\omega = \frac{-p}{\omega}
\]
folgt.

Es seien ferner zwei der Werte (3) % [sic] printed (3); the quantities under discussion are (15)
\[
\frac{\omega+c}{p}, \quad \frac{\omega+c'}{p}
\]
äquivalent, so daß eine Gleichung besteht:

\weq{26}{\frac{\omega+c}{p} = \frac{\gamma p+\delta(\omega+c')}{\alpha p+\beta(\omega+c')}, \quad \alpha\delta-\beta\gamma=1,}

oder

\weq{27}{\begin{aligned} \beta\,\omega^2 + (\beta c'+\beta c+\alpha p-\delta p)\omega \; +\phantom{0}\\ \beta c c'+\alpha c p-\delta c' p-\gamma p^2 &= 0, \end{aligned}}

woraus durch Vergleichung mit (9)

\weq{28}{\begin{aligned} \beta &= Ax,\\ \beta c'+\beta c+\alpha p-\delta p &= Bx,\\ \beta c c'+\alpha c p-\delta c' p-\gamma p^2 &= Cx, \end{aligned}}
%CONT>>
